\section{SoLM: un modelo de lenguaje basado en la fusión del AR y la difusión}

\subsection{Idea central: ordenamiento + atención causal}

La contribución fundamental de SoLM (Sorting-based Language Model, también conocido como Esoteric Language Model) es introducir soporte para el \textbf{lote} en MDLM. La idea clave es:

\begin{importantbox}{Estrategia de entrenamiento de SoLM}
Reemplazar la atención bidireccional de MDLM por \textbf{atención causal}, lográndolo mediante las siguientes operaciones:
\begin{enumerate}
    \item Realizar un \textbf{ordenamiento/barajado} del secuencia de ruido $\mathbf{z}_t$
    \item Colocar todos los \textbf{tokens limpios a la izquierda} (como contexto conocido)
    \item Colocar todos los \textbf{tokens enmascarados a la derecha} (como objetivos por predecir)
    \item Aplicar \textbf{atención causal} sobre la secuencia ordenada
\end{enumerate}
\end{importantbox}

Con esta disposición:
\begin{itemize}
    \item El primer token enmascarado depende únicamente del contexto limpio
    \item El segundo token enmascarado depende del contexto limpio + el anterior token enmascarado
    \item Y así sucesivamente...
    \item Ningún token enmascarado depende de los tokens enmascarados "futuros" a su derecha
\end{itemize}

\subsection{Proceso de inferencia y caché KV}

El proceso de inferencia de SoLM aprovecha plenamente las capacidades de la caché KV derivadas de la atención causal:

\begin{enumerate}
    \item El muestreador de antepasados decide qué posiciones enmascaradas deben ser denoiseadas (por ejemplo, las posiciones 3 y 2)
    \item Se realiza una \textbf{propagación hacia adelante únicamente para estos tokens enmascarados}, incluyendo sus codificaciones de posición correspondientes
    \item Las activaciones de los tokens ya predichos se \textbf{congelan} (se almacenan en caché)
    \item En el siguiente paso de denoising, solo es necesario procesar los nuevos tokens enmascarados, aprovechando la caché KV guardada
\end{enumerate}

\begin{importantbox}{Ahorro computacional en la inferencia de SoLM}
En comparación con MDLM, que requiere procesar toda la secuencia en cada paso (incluyendo numerosos tokens enmascarados que no transportan información), SoLM \textbf{solo procesa los tokens por denoising} en cada paso. Esto genera ahorros en dos frentes:
\begin{enumerate}
    \item La longitud de la secuencia para la propagación hacia adelante se reduce drásticamente
    \item La caché KV evita el cálculo repetido de las activaciones de los tokens ya procesados
\end{enumerate}
Ambos factores logran conjuntamente una \textbf{aceleración de la inferencia de orden de magnitud}.
\end{importantbox}

\subsection{Interpolación entre AR y difusión}

Otra contribución importante de SoLM es su capacidad para realizar una \textbf{interpolación suave} entre el modo puro AR y el modo puro de difusión:

\begin{itemize}
    \item \textbf{Modo puro AR}: en cada paso solo se denoisa un token (de izquierda a derecha), ofreciendo la mayor calidad pero con la menor velocidad
    \item \textbf{Modo puro de difusión}: en cada paso se denoisear múltiples tokens (en paralelo), logrando la máxima velocidad aunque pueda haber pérdida de calidad
    \item \textbf{Modo intermedio}: permite un equilibrio flexible entre velocidad y calidad según las necesidades
\end{itemize}

\begin{knowledgebox}{Equilibrio flexible entre velocidad y calidad}
Esta capacidad de interpolación significa que el usuario puede seleccionar el punto de operación óptimo según los requisitos del escenario de aplicación específico. Por ejemplo, en servicios en línea sensibles a la latencia se puede optar por una generación paralela más agresiva, mientras que en escenarios con exigencias extremas de calidad se adopta un modo cercano al AR puro.
\end{knowledgebox}

\subsection{Resumen de la sección}

A través de un diseño ingenioso basado en ordenamiento + atención causal, SoLM logra el soporte para caché KV manteniendo la capacidad de generación paralela de los modelos de difusión. Esto permite una mejora de la eficiencia de inferencia de orden de magnitud, haciendo que los modelos de lenguaje basados en difusión finalmente alcancen velocidades de inferencia prácticas para su uso real.

\newpage
